\subsection{Дальний восток}
Как уже было сказано Великобритания поддерживала политику "открытых дверей". Что было закреплено в Янцзинском соглашении  (англ. Yangtze Agreement) — дипломатическое соглашение между Великобританией и Германской империей, подписанное 16 октября 1900 года в разгар Боксёрского восстания. Оно закрепляло совместную позицию европейских держав о поддержании принципа «открытых дверей» и территориальной целостности Китая. А всё из-за того, что эти империи уже имели развитую промышленность и в прямой конкуренции они могли бы занять весь рынок Китая без прямого захвата региона. 

Поэтому в 1902 году был заключен Англо-Японский альянс, который обязывал войти в войну союзника в войне с двумя и более врагами. Сделано это было против России т. к. она обязывалась защищать Китай. Заручившись поддержкой, Япония в 1904 решила начать войну против Российской империи. Россия была не готово к такому конфликту: правительство не ожидало войны с Японией, считая что они не посмеют это сделать, военных сил в дальневосточной части страны не хватало для отражения атаки, основная битва проходила на море, а флот Москвы был разделен на 3 части: Балтийскую, Черноморскую (которой запрещено проходить через турецкие проливы) и Тихоокеанскую.

Одним из самых сильных кораблей в Тихом океане были выведены из строя в первые дни войны, т.к. Японцы воспользовались эффектом неожиданности и напали на Потр-Артур, где располагались корабли. У России не получилось мобилизовать силы, из-за чего война была проиграна и в 1905 году Москва утратила свои владения в Маньчжурии, на Ляодунском полуострове, а также южную часть Сахалина, и признала Корею Японской сферой влияния. После войны стало понятно, что Российская империя находится в кризисе, и сейчас не обладает лучшей силой.